\documentclass[11pt,a4paper]{article}
\usepackage[utf8]{inputenc}
\usepackage[T1]{fontenc}
\usepackage{amsmath,amssymb,amsthm}
\usepackage{geometry}
\usepackage{hyperref}
\usepackage{booktabs}
\geometry{margin=1in}
\hypersetup{colorlinks=true,linkcolor=blue,citecolor=blue,urlcolor=blue}

\newtheorem{theorem}{Theorem}
\newtheorem{conjecture}{Conjecture}
\newtheorem{proposition}{Proposition}
\newtheorem{observation}{Empirical Observation}
\newtheorem{definition}{Definition}
\newtheorem*{emplaw}{Empirical Law}

\title{The Prime Gap Dynamical Dynamical Dynamical Dynamical Dynamical Transition Operators over Prime Gap Ensembless over Prime Gap Ensembless over Prime Gap Ensembless over Prime Gap Ensembless over Prime Gap Ensembles on $(\mathbb{Z}/m\mathbb{Z})^*$:\\
Empirical Kernel Law, Compression Law, and Documented Falsifications}

\author{Bradley Wallace\\
Koba42 Research / TensorRent\\
\texttt{bwallace@tensorrent.io}}

\date{May 2026}

\begin{document}
\maketitle

\begin{abstract}
We study the empirical second-order Dynamical Dynamical Dynamical Dynamical Dynamical Transition Operators over Prime Gap Ensembless over Prime Gap Ensembless over Prime Gap Ensembless over Prime Gap Ensembless over Prime Gap Ensembles $T_m$ describing prime-gap dynamics on the multiplicative residue group $(\mathbb{Z}/m\mathbb{Z})^*$. Writing $T_m = T_0 + P_m$ where $T_0$ is the uniform transport respecting residue continuity, we establish two empirical laws verified across moduli $m$ with $\varphi(m) \leq 24$:

\textbf{Kernel Law:} $\dim \ker(P_m) = \varphi(m)$ exactly, with kernel equal to the space of state-space functions depending only on the source residue class. A structural proposition (Section~\ref{sec:structural}) establishes that the source-only sector lies in $\ker(P_m)$ as a consequence of column stochasticity and residue continuity; the empirically nontrivial content is the exact dimensional saturation across all tested moduli with no excess kernel.

\textbf{Compression Law (pre-asymptotic):} Over the tested range, the effective dynamical rank satisfies $r_{\text{eff}}(m) \sim \varphi(m)^{\beta}$ with $\beta \approx 1.6$, while the algebraic rank grows as $\varphi(m)(\varphi(m)-1)$. The ratio $r_{\text{eff}}/r_{\text{alg}}$ decays consistently with a power law of exponent $\approx -0.64$ on the tested range. The active spectrum is consistent with a heavy-tailed distribution. We make no asymptotic claim outside the tested range.

We document three conjectures tested and falsified during the investigation: uniform spectral contraction $\rho(P_m) < 1$, full character-product block-diagonalization, and universal renormalized spectral law. Finally, Section~\ref{sec:bridge} records a phenomenological closure ansatz relating the variance floor in cross-residue transition statistics to the cross-correlation bound from \cite{WallaceB}; this is presented as a numerical match, not as a derived consequence, and is quarantined as interpretive material.

Computational data covers $5.08 \times 10^7$ primes up to $10^9$ for low-modulus statistics and $1.27 \times 10^6$ primes up to $2 \times 10^7$ for the full operator-theoretic analysis.
\end{abstract}

\section{Introduction}
\label{sec:intro}

The distribution of prime gaps has been studied extensively through probabilistic models (Cram\'{e}r \cite{Cramer}, Granville \cite{Granville}) and through the explicit formula connecting primes to nontrivial zeros of the Riemann zeta function and Dirichlet $L$-functions. We approach the problem from a third angle: by constructing an explicit Markov-type Dynamical Dynamical Dynamical Dynamical Dynamical Transition Operators over Prime Gap Ensembless over Prime Gap Ensembless over Prime Gap Ensembless over Prime Gap Ensembless over Prime Gap Ensembles on the multiplicative residue group $(\mathbb{Z}/m\mathbb{Z})^*$ and identifying its algebraic and spectral structure directly from empirical prime data.

This note records the results of a sustained computational investigation of this operator. The investigation passed through three rounds of explicit falsification, each of which eliminated a conjectured structure and sharpened the remaining claims. We present here both the surviving empirical laws (the Kernel Law and the Compression Law of Sections~\ref{sec:kernel} and~\ref{sec:compression}) and the documented negative results (Section~\ref{sec:falsified}). The negative results are essential context: they constrain the theoretical interpretation of the surviving structure and rule out several natural conjectures.

The connection to the existing ACS framework \cite{Wallace2026a, WallaceB, WallaceC} is limited and quarantined to Section~\ref{sec:bridge}: a phenomenological closure ansatz that matches the empirical variance floor of cross-residue transitions on the mod-6 hexagonal clock to Paper B's cross-correlation bound on the Riemann zero spectrum. The match is striking but is not a derivation; the operator-theoretic content of the present paper is independent of this material.

\section{Setup}
\label{sec:setup}

Let $m \geq 2$ and $U_m = (\mathbb{Z}/m\mathbb{Z})^* = \{a \in \{1, \ldots, m-1\} : \gcd(a, m) = 1\}$, with $|U_m| = \varphi(m)$. Define the \emph{transition state space}
\begin{equation}
\mathcal{S}_m = \{(a, b) : a, b \in U_m\}
\end{equation}
with $|\mathcal{S}_m| = \varphi(m)^2$. For each prime gap $p_{i+1} - p_i$ with $p_i, p_{i+1} > m$ and $\gcd(p_i, m) = \gcd(p_{i+1}, m) = 1$, the pair $(p_i \bmod m, p_{i+1} \bmod m) \in \mathcal{S}_m$ is the \emph{transition class}.

\subsection{The empirical operator}

Construct the empirical second-order Dynamical Dynamical Dynamical Dynamical Dynamical Transition Operators over Prime Gap Ensembless over Prime Gap Ensembless over Prime Gap Ensembless over Prime Gap Ensembless over Prime Gap Ensembles $T_m : \mathbb{C}^{\mathcal{S}_m} \to \mathbb{C}^{\mathcal{S}_m}$ by
\begin{equation}
T_m\bigl[(a',b'), (a,b)\bigr] = \frac{N\bigl((a,b) \to (a',b')\bigr)}{\sum_{(c,d)} N\bigl((a,b) \to (c,d)\bigr)}
\end{equation}
when $b = a'$ (residue continuity), and zero otherwise. Here $N((a,b) \to (a',b'))$ counts the number of times the transition pair $(a,b)$ is followed by $(a',b')$ in the empirical sequence of consecutive transitions among primes. $T_m$ is column-stochastic on its sparsity pattern.

\subsection{The reference and the perturbation}

The \emph{unbiased reference} $T_0$ assigns uniform $1/\varphi(m)$ probability to each consistent next state:
\begin{equation}
T_0\bigl[(b,c), (a,b)\bigr] = \frac{1}{\varphi(m)} \text{ for all } c \in U_m.
\end{equation}

The empirical perturbation is
\begin{equation}
P_m = T_m - T_0.
\end{equation}
By construction, the column sums of $P_m$ are zero, and $P_m$ is supported on the same sparsity pattern as $T_m$.

\section{The Kernel Law}
\label{sec:kernel}

\begin{emplaw}[Kernel Law]
\label{thm:kernel}
For all tested moduli $m \in \{6, 10, 12, 14, 15, 18, 20, 24, 30, 42, 60, 66, 84, 90\}$:
\begin{equation}
\dim \ker(P_m) = \varphi(m).
\end{equation}
Moreover, the kernel equals the \emph{source sector}:
\begin{equation}
\ker(P_m) = \bigl\{f \in \mathbb{C}^{\mathcal{S}_m} : f(a, b) = g(a) \text{ for some } g : U_m \to \mathbb{C}\bigr\}.
\end{equation}
Dirichlet characters on $U_m$ furnish a canonical orthogonal basis for this sector, but the structural object is the source sector itself; the character basis is a coordinate choice on it.
\end{emplaw}

\subsection{Structural source invariance}
\label{sec:structural}

To delineate what is automatic from what is empirically nontrivial, we record a structural proposition.

\begin{proposition}[Structural source invariance]
\label{prop:structural}
Let $T$ and $T_0$ be column-stochastic operators on $\mathbb{C}^{\mathcal{S}_m}$ such that:
\begin{enumerate}
\item (Residue continuity) Both $T$ and $T_0$ satisfy $T[(a',b'),(a,b)] = 0$ whenever $b \neq a'$.
\item (Uniform reference) $T_0[(b,c),(a,b)] = 1/\varphi(m)$ for all $c \in U_m$.
\end{enumerate}
Then for any $f \in \mathbb{C}^{\mathcal{S}_m}$ of the form $f(a,b) = g(a)$ (source-only),
\begin{equation}
(T - T_0) f = 0
\end{equation}
in the sense that $(T-T_0)f$ is supported only on transitions inconsistent with residue continuity (which carry zero mass in both $T$ and $T_0$).
\end{proposition}

\begin{proof}[Sketch]
Both $T$ and $T_0$ are forward-transport operators on transition pairs: they map an incoming state $(a,b)$ to an outgoing state $(b,c)$. A source-only function $f(a,b) = g(a)$ assigns value $g(a)$ to the incoming state and (by extension) value $g(b)$ to the outgoing state. Column stochasticity ensures $\sum_c T[(b,c),(a,b)] = 1$, so when $(T f)$ is evaluated at outgoing state $(b,c)$, the result depends only on the source coordinate $b$ of the outgoing state. The same holds for $T_0$. The difference $(T-T_0)f$ therefore evaluates to zero on every state consistent with residue continuity.
\end{proof}

Proposition~\ref{prop:structural} establishes that the source sector lies in $\ker(P_m)$ automatically; this is a coordinate-level consequence of the construction. The empirically nontrivial content of the Kernel Law is therefore not the existence of this invariant sector, but the \emph{exact dimensional saturation}: across all tested moduli, $\dim \ker(P_m) = \varphi(m)$ with no excess kernel beyond the source sector. The empirical perturbation $P_m$ is therefore an exact destination-only perturbation in the precise sense that no additional invariants emerge from the arithmetic data.

\subsection{Verification}

The dimension match $\dim \ker(P_m) = \varphi(m)$ holds across all 14 tested moduli with no exceptions. The explicit identification of the kernel as the source sector was verified at $m = 42$ ($\varphi(m) = 12$) by constructing all 12 Dirichlet characters $\chi_1, \ldots, \chi_{12}$ on $U_{42}$ and testing source-extension vectors $v_\chi(a, b) = \chi(a)$ and destination-extension vectors $w_\chi(a, b) = \chi(b)$.

Result:
\begin{itemize}
\item $\|P_{42} v_\chi\| < 10^{-3}$ for all 12 source-extension vectors (numerical zero);
\item $\|P_{42} w_\chi\| > 0.05$ for all 12 destination-extension vectors (clearly nonzero).
\end{itemize}

This confirms that the source sector is annihilated as predicted by Proposition~\ref{prop:structural} and that no excess kernel is present at this modulus.

\subsection{The empirically nontrivial claim}

The substantive empirical content of the Kernel Law is that $P_m$ exhibits no kernel dimension beyond the structural $\varphi(m)$-dimensional source sector. Equivalently: the empirical arithmetic data, when summarized as a Dynamical Dynamical Dynamical Dynamical Dynamical Transition Operators over Prime Gap Ensembless over Prime Gap Ensembless over Prime Gap Ensembless over Prime Gap Ensembless over Prime Gap Ensembles, contains no additional invariants beyond those forced by the construction. This is a nontrivial constraint on the perturbation: there exist column-stochastic operators satisfying residue continuity for which $\dim \ker(P_m) > \varphi(m)$ (for instance, operators whose destination distributions also factor through source-coordinate symmetries). The empirical $P_m$ does not exhibit such factorizations across any tested modulus.

\section{The Compression Law}
\label{sec:compression}

\begin{definition}
The \emph{algebraic rank} of $P_m$ is $r_{\text{alg}}(m) = \mathrm{rank}(P_m) = \varphi(m)^2 - \varphi(m)$. The \emph{effective dynamical rank} (participation ratio) is
\begin{equation}
r_{\text{eff}}(m) = \frac{\bigl(\sum_i |\lambda_i|\bigr)^2}{\sum_i |\lambda_i|^2}
\end{equation}
where the sum runs over the nonzero eigenvalues of $P_m$.
\end{definition}

\begin{observation}[Empirical scaling, pre-asymptotic]
\label{obs:scaling}
Across the tested range:

\begin{center}
\begin{tabular}{cccccc}
\toprule
$m$ & $\varphi(m)$ & $r_{\text{alg}}$ & $r_{\text{eff}}$ & $r_{\text{eff}}/r_{\text{alg}}$ \\
\midrule
6 & 2 & 2 & 1.53 & 0.767 \\
12 & 4 & 12 & 7.07 & 0.589 \\
30 & 8 & 56 & 15.33 & 0.274 \\
42 & 12 & 132 & 36.52 & 0.277 \\
60 & 16 & 240 & 42.98 & 0.179 \\
90 & 24 & 552 & 100.48 & 0.182 \\
\bottomrule
\end{tabular}
\end{center}

Log-log fits over the tested range $\varphi(m) \in \{2, 4, 8, 12, 16, 24\}$ yield:
\begin{align}
r_{\text{eff}}(m) &\approx 0.59 \cdot \varphi(m)^{\beta}, \quad \beta \approx 1.61,\\
\frac{r_{\text{eff}}(m)}{r_{\text{alg}}(m)} &\approx 1.23 \cdot \varphi(m)^{-\delta}, \quad \delta \approx 0.64.
\end{align}
\end{observation}

\textbf{Scope.} These exponents are pre-asymptotic fits over a small modulus range and should not be extrapolated to large $\varphi(m)$ without further verification. Within the tested range, the ratio $r_{\text{eff}}/r_{\text{alg}}$ is monotonically decreasing, and we report this as an empirical observation only. Whether the trend continues, saturates, or reverses for $\varphi(m) > 24$ is an open computational question.

The qualitative content of the observation is that the active sector exhibits a substantial compression: while the algebraic rank of $P_m$ grows quadratically with $\varphi(m)$, only a sub-quadratic number of modes carries the dynamics over the tested range. The asymptotic regime is not established by the present data.

\section{Spectral Concentration}
\label{sec:spectrum}

The distribution of nonzero eigenvalue magnitudes $|\lambda|$ of $P_m$ is concentrated on a small number of dominant modes. Combining data from $m \in \{42, 60, 90\}$ yields 924 nonzero eigenvalues. The quantiles are summarized below.

\begin{center}
\begin{tabular}{cc}
\toprule
Quantile & $|\lambda|$ \\
\midrule
50\% & 0.015 \\
75\% & 0.025 \\
90\% & 0.055 \\
95\% & 0.124 \\
99\% & 0.323 \\
\bottomrule
\end{tabular}
\end{center}

The 99th percentile is approximately $20\times$ the median; the top 1\% of eigenvalues span a range larger than the entire bulk distribution. A small number of dominant modes carry most of the dynamical mass.

A pure power-law fit on the top 30\% of eigenvalues gives a tail exponent in the neighborhood of $-1.2$. We do not present this as an established law: the sample is small, alternative parametric families (lognormal, stretched exponential) have not been distinguished, and goodness-of-fit testing has not been performed. The qualitative observation is that the distribution is consistent with a heavy tail and dominated by a small number of large modes; the precise functional form is open.

This concentration pattern is consistent with the low-modulus anomalies on the mod-6 hexagonal clock that motivated the study. At $\varphi(m) = 2$, the active sector contains only one or two modes, all of them dominant; the strong empirical statistics at low modulus are the signature of these few modes operating with little background. At larger modulus, the dominant modes persist but are joined by a high-dimensional bulk of weakly active modes.

\section{Phenomenological Closure Ansatz and the ACS Cross-Reference}
\label{sec:bridge}

This section records a numerical coincidence between the empirical variance floor of cross-residue transitions on the mod-6 hexagonal clock and the cross-correlation bound from \cite{WallaceB}. We present it as a phenomenological closure ansatz, not as a derived consequence. The material is interpretive and is included here to record the empirical match; readers concerned only with the operator-theoretic content of the paper may skip this section without loss.

\textbf{The empirical numbers.} The investigation began with three statistical anomalies on the mod-6 hexagonal clock at $5 \times 10^7$ primes: a sign-reversed restoring force (gaps systematically shorten after multi-rotation events), a directional asymmetry (the $5 \to 1$ transition is 12.75\% faster than $1 \to 5$ at $10^7$, attenuating to 7.7\% at $10^9$), and a variance suppression of cross-class transitions to $0.292 \pm 0.003$ relative to the value expected under independent sampling (measured at window size $W = 100$ on the $10^9$ Rust sieve data, stable across $W \in [100, 500]$).

\textbf{The ansatz.} For partially correlated indicators with pairwise correlation magnitude $C$ and effective number $N_{\text{eff}}$ of correlated pairs, a standard mean-field closure relates the variance of the empirical fraction to:
\begin{equation}
R = \frac{1}{1 + (N_{\text{eff}} - 1) \cdot |C|}.
\label{eq:closure}
\end{equation}
This is a phenomenological relation: it is not derived from the explicit formula or any rigorous spectral identity. It is a standard mean-field closure that arises in many partially-correlated-sampling contexts.

\textbf{The numerical match.} Substituting the empirical autocorrelation-derived value $N_{\text{eff}} = 56$ (extracted from the cross-class indicator autocorrelation on prime data, $W = 100$) and Paper B's cross-correlation bound $|C_N| = 0.044$ from the first 200 Riemann zeros \cite{WallaceB}:
\begin{equation}
R = \frac{1}{1 + 55 \times 0.044} = \frac{1}{3.42} = 0.292.
\end{equation}
This matches the empirical variance ratio to three significant figures.

\textbf{What this is and what it is not.} The match is striking but is not a derivation. The closure formula~\eqref{eq:closure} is not derived from the Wronskian non-vanishing of \cite{WallaceB}; the identification of $N_{\text{eff}} = 56$ with a quantity computed from the Riemann zero spectrum is not established; and the appearance of $|C_N| = 0.044$ on both sides of the match is suggestive but not causal. A rigorous derivation would require deriving~\eqref{eq:closure} from the explicit formula and showing that the relevant $N_{\text{eff}}$ for the prime-gap autocorrelation in fact coincides with the spectral cross-correlation parameter. Neither step is performed in this paper.

The cross-reference is included as a quantitative coincidence worth recording; it is not load-bearing for the Kernel Law or Compression Law, which stand independently as empirical results about $P_m$.

The other two anomalies (sign-reversed spring-back, directional asymmetry) admit mechanistic identification with Dirichlet character corrections and zero-mode decoherence, respectively, but lack the same level of numerical match. These remain open problems (Section~\ref{sec:open}).

\section{Three Falsified Conjectures}
\label{sec:falsified}

The investigation proposed and explicitly tested three conjectures that did not survive the data. We document them because their failure modes constrain the theoretical interpretation of the surviving structure.

\subsection{Uniform spectral contraction (falsified)}

\textbf{Conjecture:} $\rho(P_m) < 1$ uniformly in $m$, with $P_m$ behaving as a quasi-nilpotent contraction.

\textbf{Falsification:}
\begin{center}
\begin{tabular}{ccc}
\toprule
$m$ & $\varphi(m)$ & $\rho(P_m)$ \\
\midrule
6 & 2 & 0.145 \\
30 & 8 & 0.341 \\
60 & 16 & 0.586 \\
90 & 24 & 0.745 \\
\bottomrule
\end{tabular}
\end{center}
Power-law fit: $\rho(P_m) \approx 0.07 \cdot \varphi(m)^{0.75}$, predicting $\rho = 1$ near $\varphi(m) \approx 50$. The contraction principle holds in the low-modulus regime but does not extrapolate.

\subsection{Character-product block-diagonalization (falsified)}

\textbf{Conjecture:} $P_m$ acts block-diagonally in the basis $\{\chi_i(a) \overline{\chi_j(b)}\}$ of $\varphi(m)^2$ character products on the state space.

\textbf{Falsification:} Define the \emph{left-preserved norm fraction} as the ratio of $\|P_m\|_F$ contributions from entries respecting left-character labels to the total $\|P_m\|_F$.

\begin{center}
\begin{tabular}{cc}
\toprule
$m$ & Left-preserved fraction \\
\midrule
6 & 71\% \\
12 & 53\% \\
30 & 36\% \\
42 & 29\% \\
\bottomrule
\end{tabular}
\end{center}
The fraction \emph{decreases} with $\varphi(m)$. At $m = 42$, only 29\% of the perturbation norm respects left-character blocks. Character structure is exact at the kernel boundary (the Kernel Law) but does not extend to global block-diagonalization.

\subsection{Universal renormalized spectral law (falsified)}

\textbf{Conjecture:} The normalized operator $P_m / \rho(P_m)$ has a universal spectral distribution as $m \to \infty$.

\textbf{Falsification:} The mean of $|\lambda|/\rho(P_m)$ over the nonzero spectrum:
\begin{center}
\begin{tabular}{ccc}
\toprule
$m$ & $\varphi(m)$ & Mean $|\lambda|/\rho$ \\
\midrule
6 & 2 & 0.645 \\
12 & 4 & 0.402 \\
30 & 8 & 0.126 \\
60 & 16 & 0.054 \\
90 & 24 & 0.040 \\
\bottomrule
\end{tabular}
\end{center}
The normalized mean decreases monotonically by a factor of 16. No universal renormalized law exists in this normalization. The spectrum becomes more concentrated on dominant modes as $\varphi(m)$ grows.

\section{Open Problems}
\label{sec:open}

\subsection{Result hierarchy}

To make the epistemic status of the various claims in this paper explicit, we record the tier classification used throughout:

\begin{center}
\begin{tabular}{p{0.42\textwidth}p{0.5\textwidth}}
\toprule
\textbf{Result} & \textbf{Status} \\
\midrule
$\dim \ker(P_m) = \varphi(m)$ (Kernel Law) & Verified across 14 moduli; Proposition~\ref{prop:structural} explains source-sector inclusion; the empirical saturation is non-trivial \\
Source sector identified as kernel & Verified explicitly at $m = 42$, exhausting Dirichlet characters \\
$r_{\text{eff}}(m) \sim \varphi(m)^{1.6}$ scaling & Empirical fit on tested range; not extrapolated \\
Heavy-tail spectrum, exponent near $-1.2$ & Empirical observation; precise tail form open \\
ACS variance-floor bridge (Section~\ref{sec:bridge}) & Phenomenological closure, numerical match only \\
Uniform contraction $\rho(P_m) < 1$ & Falsified --- $\rho(P_m)$ grows with $\varphi(m)$ \\
Character-product block-diagonalization & Falsified --- character separation holds only at kernel level \\
Universal renormalized spectral law & Falsified --- distribution concentrates with $\varphi(m)$ \\
\bottomrule
\end{tabular}
\end{center}

\subsection{The structural mystery}

The deepest open question raised by the data is not the spectral radius behavior or the compression exponent. It is this:

\begin{quote}
\emph{Why does the operator $P_m$ preserve source-character invariants exactly while destroying character separability everywhere else?}
\end{quote}

The kernel is the source-character span, exactly. But on the orthogonal complement, the perturbation does \emph{not} respect the character-product basis $\{\chi_i(a) \overline{\chi_j(b)}\}$ --- as documented in Section~\ref{sec:falsified}, the left-preserved norm fraction drops to 29\% at $m = 42$. The operator therefore exhibits an asymmetric harmonic structure: complete invariance at the boundary of the kernel, near-complete mixing in the interior.

This suggests $P_m$ may be best understood as a non-normal operator with arithmetic harmonic boundary conditions: a conditional-expectation-like structure on the source coordinate combined with non-semisimple mixing on the destination coordinate. Finding the right mathematical category --- whether this is a Markov dilation, an arithmetic Toeplitz structure, or something more exotic --- is the most promising direction for converting the empirical Kernel Law into a theorem.

\subsection{Specific computational targets}

\textbf{1. Derive the Kernel Law.} The empirical exactness of $\dim \ker(P_m) = \varphi(m)$ across all tested moduli, beyond Proposition~\ref{prop:structural}'s inclusion result, strongly suggests an algebraic theorem. A rigorous derivation would establish that no additional invariants arise from the arithmetic structure of the empirical transition data.

\textbf{2. Constrain the Compression Law.} Whether $r_{\text{eff}} \sim \varphi(m)^{1.6}$ is a fundamental scaling or an intermediate asymptotic regime is open. Testing $\varphi(m) > 100$ would constrain extrapolation.

\textbf{3. Connection to $L$-functions.} The dominant low-modulus mode at $m = 6$ has spectral radius 0.145, attenuation rate consistent with $1/\log(x)$ scaling, and is structurally compatible with the $L(s, \chi_3)$ explicit formula contribution to the Chebyshev bias. A rigorous derivation of mode eigenvalues from $L$-function zero spectra remains open.

\textbf{4. Asymptotic spectral law.} The tail exponent near $-1.2$ should be measured at larger $m$ to distinguish a fixed limiting law from a slowly varying one.

\textbf{5. Spring-back derivation.} The sign-reversed restoring force (gaps systematically shorten after multi-rotation events) is consistent with zero-mode decoherence --- the inability of the spectral components to maintain destructive alignment --- but lacks a quantitative bridge analogous to the variance floor. Its derivation likely requires local Gaussian Unitary Ensemble (GUE) zero-spacing statistics integrated with the global spectral structure.

\section{Honest Scope Boundary}

This note establishes empirical regularities of an operator constructed from prime-gap data. It does not constitute a proof of any number-theoretic theorem about primes. The Kernel Law is verified across 14 moduli and explicitly mechanism-tested at $m = 42$; it is presented as an empirical theorem awaiting rigorous proof. The Compression Law is a fitted scaling law; the underlying mechanism is open.

The conjectures that did not survive — uniform contraction, character-product block-diagonalization, and universal renormalized law — were tested and explicitly rejected by the data. Their failure modes are documented because negative results constrain the theoretical interpretation of the positive results.

\section{Data and Reproducibility}

All computations use primes generated by the Sieve of Eratosthenes implemented in Python (up to $2 \times 10^7$) and Rust (up to $10^9$ for the low-modulus statistics). The Rust sieve is segmented and parallelized via Rayon; the Python implementation is single-threaded and suffices for the operator-theoretic range.

The operator construction, kernel verification, and scaling-law extraction are deterministic and reproducible from the source primes.

\section{Summary}

The empirical prime-gap Dynamical Dynamical Dynamical Dynamical Dynamical Transition Operators over Prime Gap Ensembless over Prime Gap Ensembless over Prime Gap Ensembless over Prime Gap Ensembless over Prime Gap Ensembles $P_m = T_m - T_0$ on the multiplicative residue group $(\mathbb{Z}/m\mathbb{Z})^*$ exhibits the following structure across moduli with $\varphi(m) \leq 24$:

\emph{Kernel Law (verified across 14 moduli):} The kernel of $P_m$ is exactly $\varphi(m)$-dimensional and consists of source-residue-only functions. Proposition~\ref{prop:structural} accounts for the inclusion of the source sector in the kernel; the empirically nontrivial content is the exact dimensional saturation across all tested moduli, with no excess kernel beyond the source sector.

\emph{Compression Observation (pre-asymptotic, tested range):} The effective dynamical rank satisfies $r_{\text{eff}}(m) \sim \varphi(m)^{\beta}$ with $\beta \approx 1.6$ over the fitted range; the algebraic rank grows as $\varphi(m)(\varphi(m) - 1)$. The active spectrum is concentrated on a small number of dominant modes, with the 99th percentile of $|\lambda|$ exceeding $20 \times$ the median. The precise functional form of the tail and its asymptotic behavior are open.

\emph{Documented falsifications:} Uniform spectral contraction (the data instead shows $\rho(P_m)$ growing approximately as $\varphi(m)^{0.75}$), character-product block-diagonalization (only kernel-level character invariance survives), and universal renormalized spectral law (the renormalized distribution becomes progressively concentrated as $\varphi(m)$ grows).

\emph{Phenomenological ansatz (Section~\ref{sec:bridge}, quarantined):} A mean-field closure relation matches the empirical variance floor of cross-class transitions to the cross-correlation bound from \cite{WallaceB}. This is recorded as a numerical match worth investigating, not as a derived relation.

The structural content of the paper is the Kernel Law combined with Proposition~\ref{prop:structural}; the empirical scaling and concentration phenomena complement this with information about how $P_m$ behaves outside its kernel.

\begin{thebibliography}{99}

\bibitem{Wallace2026a}
B.~Wallace,
\emph{Colour from Gravity: Pati-Salam from the Palatini Bracket},
TensorRent, 2026.

\bibitem{WallaceB}
B.~Wallace,
\emph{Spectral Susceptibility and Renormalized Stability on the Riemann Critical Line},
TensorRent, 2026.

\bibitem{WallaceC}
B.~Wallace,
\emph{Holographic Spectral Inversion and Invariant Kinematic Attractors},
TensorRent, 2026.

\bibitem{Cramer}
H.~Cram\'{e}r,
\emph{On the order of magnitude of the difference between consecutive prime numbers},
Acta Arith. 2 (1936), 23--46.

\bibitem{Granville}
A.~Granville,
\emph{Harald Cram\'{e}r and the distribution of prime numbers},
Scand.\ Actuar.\ J.\ 1995, no.~1, 12--28.

\end{thebibliography}

\end{document}
